\section{Introduction}
\label{sec:intro}

Large language models (LLMs) are now capable of writing programs that recompute processes with changes applies to answer what-if questions. that does this recomputing. Given a document that states a process's rules and records its history, a capable model can produce a simulator: code that reads the recorded inputs, applies the rules, and reports what happens~\cite{otftools2026,chen2026devs,holt2025gsim}. Once built, the same program can answer many questions, including ones about new records or a revised policy. Building it once therefore costs far less than having a model reason through the process for each question~\cite{otftools2026}.

This paper studies whether these reusable programs can be trusted. Before a generated program is used, platforms check it. A common check, and a classic test of simulation models~\cite{sargent2013vv}, is whether the program reproduces the recorded history: hide some recorded outcomes, run the program, and compare. We call this a \emph{history check}. A program that passes is accepted and then used for every later question. That verdict is about the program as a whole. A what-if question, however, asks about one particular change, and a program can reproduce history perfectly while handling that change wrongly.

Consider a loan-application process from our benchmark, in which every applicant takes exactly one day to complete a form. During development, a program written by one of our models reproduced every recorded outcome on all four development runs, and the platform's own check accepted it with a perfect score. Reading its code showed why it could not be trusted for every question: it ignored the per-case durations it read from the records and used the published one-day constant instead. Because that duration never varies in the history, no history check can see the problem. Asked ``what if this applicant had finished the form in one second?'', the program answers a full day too late, and the error spreads to other cases that wait for the same workers.

Machine learning has a name for this situation: models that fit the same data equally well can still behave differently on inputs the data never exercised~\cite{damour2022underspecification}. A single pass-or-fail verdict on a program therefore says little about whether a \emph{particular answer}, for a \emph{particular change}, can be trusted. Failures can come from several places: the question may be misread, the change may be applied to the wrong record, the program may ignore the changed input or simulate its effects incorrectly, or the program may stop being valid when the policy it was built for changes. A checker that only judges the program as a whole cannot tell these apart, and cannot tell a user which answers to rely on.

\paragraph{Our approach.} We built Vouch, an assurance library that decides, for each answer, whether to trust it or decline it, based on the change the question asks for. Because the change is an intervention on the process, we call this intervention-level assurance. For each question, Vouch:
\begin{enumerate}[leftmargin=*,nosep]
  \item turns the question into a structured change (what is changed, how, by how much, when, and for which record) and checks that change against the document and the current records;
  \item checks that the program reproduces the recorded history of the run being asked about;
  \item runs small tests aimed at the requested \emph{kind} of change, built only from the document's stated rules and records, never from reference answers; for example, it changes the same kind of input for another case and checks that the program's output moves, and that nothing before the change moves;
  \item keeps a record of which kinds of change each program has passed or failed, so that a failure blocks later answers that depend on the same kind of change, and repeats the relevant tests when the program is reused on new records or under a new policy.
\end{enumerate}
Every decline comes with a reason in plain terms: a test showed the program getting this kind of change wrong (\Caught), it had failed the same test on an earlier question (\Already), no test of this kind of change could run on the program (\Untestable), the program does not reproduce the recorded history (\Failhist), or it has no input that can represent the change (\Cantexpress).

\paragraph{Evaluation.} We evaluate Vouch as one of several \emph{trust judges}, procedures that decide whether to trust an answer, all deciding about the same fixed programs, questions and answers. The other judges are each platform's own acceptance check, the history check by itself, generic tests that are not aimed at the requested change, agreement with an answer an LLM gives by reading the document, an LLM asked to judge the answer, an LLM asked to judge the program, and a ReAct agent~\cite{yao2023react} that can run the program in a sandbox before deciding. Programs come from three platforms (\otf~\cite{otftools2026}, \devs~\cite{chen2026devs}, and a plain coding agent), built by five LLMs, in four domains: manufacturing, IT operations, and loan and pharmacy workflows. Reference answers come from separately implemented and independently checked reference simulators, never from an LLM. Because a judge can always look safer by answering less, we compare judges at the \emph{same answer coverage}, and we report the cost of each.

\paragraph{Findings.} With no added check, \rNoCheckWrongClaude{} of the answers given by programs that Claude models wrote are wrong (\rNoCheckErrClaude\%). The history check leaves \rHistWrongClaude{} of them, and Vouch leaves \rGateWrongClaude{} while trusting \rGateCovClaude\% of all questions. Almost all of the difference comes from one platform, where \rOtfHistWrong{} wrong answers pass the history check and Vouch trusts none of them, at the cost of 123 of the 2{,}052 correct answers the history check keeps. Every one of these wrong answers came from a correctly interpreted question, so the faults lie in the programs. Where the history check already finds the errors, or programs make none, Vouch adds nothing. At equal coverage it accepts fewer wrong answers than LLM judges built on Claude Haiku, and ties with stronger and open-weight judges that we ran on small samples. When the rules change, here by adding one worker to a pool, every program fails the history check and Vouch answers nothing, while agreement with a reading answer keeps nearly half of the answers without error. Vouch needs no extra model calls and about \rGateExecPerQ{} extra program runs per question.

\paragraph{Contributions.}
\begin{itemize}[leftmargin=*,nosep]
  \item Vouch, a method for deciding whether to release a generated simulator's answer to a what-if question: it checks the structured request against the records, tests the requested kind of change without reference answers, consults the failures recorded for the program, and names the reason for every decline (\S\ref{sec:method}).
  \item A benchmark and protocol for answer-level assurance: 547 test questions in four domains with exact answers under fixed random draws, independently checked reference answers, no-effect controls and policy variants, and 86 programs from three platforms and five builder models, judged with the candidate answers held fixed (\S\ref{sec:problem}, \S\ref{sec:benchmark}, \S\ref{sec:judges}).
  \item An account of when answer-level checking helps and when it does not, against rule-based and LLM judges at equal coverage, reported per platform, domain, kind of change and program (\S\ref{sec:results}).
\end{itemize}
